\section{Introduction}

Small and medium-sized enterprises are economically important in countries worldwide. In Europe, for instance, SMEs represent nearly 99.8\% of all registered companies, accounting for more than half of the European GDP and two-thirds of all jobs in the private sector \citep{pwc2014}. Similar figures can be observed in Brazil, where approximately 97\% of all firms are SMEs, representing approximately 50\% of Brazil’s formal jobs \citep{bastos2018}.

In recent years, governments worldwide have implemented public policies that favor SMEs based on the potential or actual benefits that these policies can bring to the economy.

The literature widely acknowledges that SMEs have massive potential for job creation, local development, and innovation \citep{sigma2016}. However, there is limited evidence on the social costs of such policies, since most research and case studies do not address this point. This paper exploits a quasi-experimental variation from a public SME-related program to estimate the costs of incentivizing restricted public tenders to SMEs in Sao Paulo, Brazil.

The use of public procurement as a policy tool has been a major trend worldwide \citep{thai2017}. Policies that favor sustainable or `green' procurement or utilize social criteria to restrict bids to a target group of sellers are prime examples of public procurement policy that aims to achieve social and economic outcomes.

One of the most widespread practices to promote local development through public procurement is the restriction of public tenders to SMEs. The significant presence of SMEs in the economy suggests that these companies are an essential channel through which to infuse economic development \citep{mel2008,freedman2013}, although there is evidence that the effectiveness of this mechanism depends on the quality of management practices \citep{mckenzie2015} or how capital is provided through investment in SMEs \citep{fafchamps2014}.

The primary justifications for implementing policies that favor SMEs in public procurement are related to the various entry barriers to public tender that SMEs face \citep{oecd2018,loader2015,hoekman2020}. Better access to the public procurement market might catalyze SMEs’ productive potential, especially in contexts with demand constraints \citep{cardoza2016,ferraz2016}. It may expand government networks of goods and services providers, thus enhancing the competition among firms and enabling public acquisitions at more competitive prices \citep{loader2013}.

Despite the potential positive effects generated by favoring SMEs in public tenders, such a policy may undermine the quality and efficiency of the public procurement process \citep{nakabayashi2013}. First, it may lead to smaller-scale purchases per bid because SMEs generally cannot handle large orders.

Second, restricting public bid participation may harm firms’ screening process and increase the likelihood of selecting fewer, less efficient firms or no firms at all \citep{nakabayashi2013,loader2013}. Thus, either negotiated prices may be higher than usual or there may be a waste of public resources since planning and executing a bid is costly. However, there is evidence that the costs of negotiating and executing public contracts may be mitigated by public managers’ contract management capabilities and private sellers’ contract execution capabilities \citep{cabral2017}.

Additionally, depending on the characteristics of their sector, SMEs may not reach competitive price-cost levels when providing a good or service to the government \citep{nakabayashi2013}. Thus, organizing an unsuccessful tender means wasting resources; the government may incur severe costs with few results for SMEs or the local economy.

While there are numerous examples worldwide of policies favoring SMEs in public tenders, there is little direct evidence of the impacts of such policies on firms performance. Additionally, the costs of implementing such a policy are practically ignored in the literature.

Estimating the public costs of favoring SMEs in public procurement is a complicated issue to assess empirically because doing so requires establishing an appropriate counterfactual \citep{hoekman2020}. I utilize a policy experiment in the state of Sao Paulo, Brazil, to estimate the costs of favoring SMEs in public procurement. I exploit the timing of a change in a policy of restricted SME tenders (March 2018) that affect only a specific group of items (group 65). The identification strategy simultaneously uses time and cross-sectional variations to estimate the effects of this policy shift.

However, in this paper, examining how institutional change occurred allows the costs of the SME policy to be assessed only indirectly. Instead of using a standard difference-in-differences (DiD) method, I use a variation of this method known as difference-in-differences in reverse (DiDiR), or `time-reversed DiD' \citep{kim2019}. In the DiD method, there is a control group that is never treated and a treatment group that is treated at some point in time. In the DiDiR method I employ, the control group is always treated (instead of always untreated), and the other group is the `switched group,' subject to the change in policy.

In Sao Paulo, the government can execute public tenders for SMEs only. Between August 2014 and February 2018, the government’s default choice consisted of executing SME-only tenders for all items with a value less than or equal to R\$80{,}000, except for a group of items identified as code 65, which was only allowed to execute open tenders in this period.

The government can avoid restricted tenders for SMEs if it considers that any item in a bidding process falls within the exceptions provided for by law through a costly process of justification to its watchdogs. In case of any irregularity or failure to comply strictly with the law, public officers can be punished. From March 2018 on, after a change in the law’s interpretation, group 65 became subject to the general rule just as any other group of purchased items. Thus, the `always treated' group here consists of all groups of items but group 65, comprising the switched group.

Notably, DiDiR identifies pre-switch-period effects; i.e., it estimates effects for the past \citep{kim2019}, assessing the costs of the policy switch indirectly. I estimate this pre-switch-period effect on the switched group, comparing the observed outcomes for group 65 before the shift in policy and the outcomes that would have occurred for this group if there had been opt-out costs before March 2018.

Thus, the DiDiR method reveals that the negotiated prices are, on average, between 4.58\% and 8.08\% lower for group 65 than other groups before March 2008. Moreover, in the pre-period, the number of participants in group 65 is approximately 22\% higher than that in other groups of items.

The number of valid bids follows the same pattern as the number of participant firms: there are approximately 25\% more valid bids in group 65 than in the other groups of items. Finally, before the policy change, sellers that won tenders for items in group 65 were approximately 4 km further away from public buyer units (PBUs) than in open tenders.

This paper has five sections, including this introduction. Section 2 provides a brief overview of the institutional background related to SME law and public procurement in Sao Paulo, Brazil, which is relevant to the empirical section. Section 3 describes the datasets and sample definitions. Section 4 presents the empirical analysis. Section 5 concludes the paper.